% From 0% to Tested: Unit-Testing Flutter Apps with Heavy Platform Integration
% Target: IEEE Software (rolling submissions, ~4000-6000 words)
% Review: single-anonymous by default
% Format: IEEEtran journal/compsoc

\documentclass[10pt,journal,compsoc]{IEEEtran}

\usepackage{booktabs}
\usepackage{listings}
\usepackage{xcolor}
\usepackage{url}
\usepackage{cite}

\lstset{
  basicstyle=\ttfamily\small,
  breaklines=true,
  frame=single,
  language=Java,
  commentstyle=\color{gray},
  keywordstyle=\color{blue},
  numbers=left,
  numberstyle=\tiny\color{gray},
  xleftmargin=2em,
}

\begin{document}

\title{From 0\% to Tested: Unit-Testing Flutter Apps\\ with Heavy Platform Integration}

\author{Sagar~Ghag
\IEEEcompsocitemizethanks{\IEEEcompsocthanksitem S. Ghag is with Kryonex Labs, Pune, India.\protect\\
E-mail: sagar.ghag1997@gmail.com}}

\IEEEtitleabstractindextext{%

% ============================================================================
% ABSTRACT
% ============================================================================

\begin{abstract}
Flutter apps that depend on Firebase, background services, and native
platform channels present a testing problem that framework documentation
does not address: a significant fraction of production code is structurally
untestable at the unit level, and naive line-coverage targets obscure this
reality rather than measuring quality.

We report on the testing journey of Zontra, a 104,886-line production
Flutter app for field workforce tracking, from zero automated tests to
2,583 tests across 292 files with 55.6\% filtered line coverage. Our
approach starts with a coverage-exclusion policy that separates
platform-wrapper code from testable business logic. Mock boundaries using
\texttt{mocktail} and \texttt{bloc\_test} confine platform dependencies
to thin adapter layers. Seventeen custom lint rules (built on
\texttt{custom\_lint}) keep the architecture testable. We document six
cases where automated tests caught production regressions, including one
that produced a new static-analysis rule and one that exposed a
security-relevant routing bypass. We report costs: 53,354 lines of test
code, a 0.51 test-to-source line ratio, and the effect of
\texttt{fakeAsync} conversion on individual test files.

Practitioners building Flutter apps with Firebase or native integrations
can adopt the exclusion framework, mock-boundary patterns, and custom-lint
rules directly.
\end{abstract}

\begin{IEEEkeywords}
Flutter, mobile testing, platform integration, coverage policy, custom lint rules, experience report.
\end{IEEEkeywords}}

\maketitle

% ============================================================================
% 1. INTRODUCTION
% ============================================================================

\section{Introduction}
\label{sec:intro}

The Flutter testing documentation~\cite{flutter_testing_docs} presents a
clean three-layer pyramid: unit tests, widget tests, integration tests.
The examples use stateless widgets with no dependencies, or simple state
objects that accept injected collaborators. Production Flutter apps look
nothing like this.

A production app initialises Firebase before the first frame, opens
foreground services to survive background kills, registers push
notification handlers, reads from SharedPreferences, calls platform
channels into Kotlin and Swift. None of these can be instantiated in a
\texttt{flutter\_test} harness without substantial scaffolding that the
official documentation does not describe.

The result is predictable. Kochhar et al.~\cite{kochhar2015understanding}
found that mobile developers write fewer tests than their web and desktop
counterparts. Linares-V\'{a}squez et al.~\cite{linares2017developers}
confirmed that Android apps have lower test density than server-side
projects. Our experience suggests the gap is even wider for cross-platform
apps with heavy platform integration, because the framework boundary adds
a second layer of indirection that must be mocked.

This paper reports on our experience testing Zontra, a field workforce
tracking app built in Flutter with Firebase Authentication, Firestore,
Remote Config, FCM push notifications, a Rust-based background location
engine accessed via FFI, and SharedPreferences-backed local storage. The
app serves two client organisations in the Indian delivery and logistics
sector.

We started with zero tests. Eighteen months later the codebase has 2,583
tests across 292 files, covering 55.6\% of filtered source lines. The
journey was not a straight line; we over-tested in some areas, excluded
too little in the early months, and learned by watching tests break in
ways that taught us what was worth testing and what was not.

Our contributions are:

\begin{enumerate}
  \item A \textbf{coverage-exclusion framework} that separates
        platform-wrapper code from testable business logic, with a
        concrete taxonomy of what to exclude and a reporting method that
        makes the gap honest rather than hidden.
  \item A \textbf{mock-boundary strategy} for Flutter apps with Firebase
        and native platform dependencies, showing how to confine
        untestable code to thin adapters while keeping the business logic
        fully injectable.
  \item A \textbf{custom-lint package} of 17 rules that keep the
        architecture testable, including one rule born from a production
        regression.
  \item \textbf{Six documented cases} where automated tests caught real
        production bugs, including a security-relevant routing bypass.
\end{enumerate}

% ============================================================================
% 2. CONTEXT
% ============================================================================

\section{Context: The Zontra App}
\label{sec:context}

Zontra is a B2B field workforce tracking application. Employers use it to
monitor employee location during work shifts, track distance travelled,
and manage attendance via geofenced office check-ins. The app runs on
Android devices ranging from budget handsets (Redmi 9A, 3GB RAM) to
mid-range phones (Samsung Galaxy A14).

At the time of writing, the production codebase contains 750 source files
and 104,886 lines of Dart in \texttt{lib/}. Key dependencies that create
testing friction:

\begin{itemize}
  \item \textbf{Firebase}: Authentication (phone OTP), Firestore (real-time
        sync), Remote Config (feature flags), Cloud Messaging (push
        notifications). Each requires \texttt{Firebase.initializeApp()},
        which fails in a test harness without
        \texttt{firebase\_core\_platform\_interface} stubs.
  \item \textbf{Tracelet}: A Rust-based background location engine compiled
        via Flutter-Rust Bridge. The engine runs as an Android foreground
        service. It cannot be instantiated in a desktop test runner.
  \item \textbf{SharedPreferences}: Used for offline waypoint storage and
        session state. Testable only via Flutter's built-in
        \texttt{setMockInitialValues} helper, which requires explicit
        setup per test.
  \item \textbf{GetIt}: Service locator for dependency injection. All
        production dependencies are registered at startup. Tests must
        reset and re-register with mocks per test or group.
\end{itemize}

The development team is two engineers. There is no dedicated QA role.

% ============================================================================
% 3. THE TESTING WALL
% ============================================================================

\section{The Testing Wall}
\label{sec:wall}

We use ``testing wall'' to describe the point at which a developer writing
their first test for a platform-integrated Flutter app hits a cascade of
initialisation failures. The first \texttt{flutter test} invocation on our
codebase (before any test infrastructure) produced:

\begin{lstlisting}[language={},numbers=none]
FirebaseException ([core/no-app])
No Firebase App has been created
\end{lstlisting}

This error appeared not because the test exercised Firebase, but because
the file under test imported a module that transitively imported
\texttt{firebase\_core}. The dependency graph is the problem, not the test.

We identified four categories of untestable code:

\begin{enumerate}
  \item \textbf{SDK wrappers}: Thin classes that call Firebase Auth,
        Firestore, or Remote Config. These require real SDK initialisation.
        Example: \texttt{feature\_flag\_service.dart} (116 lines), which
        wraps \texttt{FirebaseRemoteConfig} fetch-on-start and a
        real-time listener stream. No unit test exists for this file.
  \item \textbf{Native bridges}: The Tracelet background location engine
        (FFI to Rust). Cannot run on a desktop test runner. Stubbed to a
        one-liner no-op: \texttt{Future<void> noOpForceStop() async \{\}}.
  \item \textbf{DI wiring}: GetIt registration files
        (\texttt{service\_locator.dart}, 455 lines across 3 files). These
        are lists of \texttt{registerLazySingleton} calls. No branching
        logic; testing them would mean testing GetIt itself.
  \item \textbf{Route tables}: \texttt{app\_router.dart} (253 lines).
        Static mapping from route names to screen constructors. The app
        originally used Fluro and later migrated to GoRouter; both
        routers are declarative maps. Testing route resolution is the
        router library's job.
\end{enumerate}

Together these categories account for 21 files and 4,154 lines, roughly
13\% of the instrumented source. The bulk of the exclusion (2,939 lines)
is the location engine's Dart-side FFI bindings: not thin wrappers, but
the full bridge layer between the app and the Rust core. The transitive
import graph means that without explicit architectural boundaries, these
files poison the test harness for everything that depends on them.

% ============================================================================
% 4. APPROACH
% ============================================================================

\section{Approach}
\label{sec:approach}

We separated what to exclude from what to test, drew mock boundaries
around the platform dependencies that remained, and wrote static-analysis
rules to keep both decisions from eroding.

% ----------------------------------------------------------------------------
\subsection{Coverage Exclusion Policy}
\label{sec:exclusion}

We run \texttt{flutter test --coverage} in CI and post-process the
resulting \texttt{lcov.info} with \texttt{lcov --remove} before computing
metrics. The exclusion patterns are declared in the project
\texttt{Makefile}, not hidden in CI scripts:

\begin{table}[h]
\caption{Coverage exclusion patterns}
\label{tab:exclusions}
\small
\begin{tabular}{@{}p{4.2cm}l@{}}
\toprule
Pattern & Reason \\
\midrule
\texttt{*.g.dart}          & Generated code \\
\texttt{firebase\_options*} & Build-generated config \\
\texttt{firebase/*.dart}   & SDK wrappers (3 files) \\
\texttt{feature\_flag\_service} & Remote Config wrapper \\
\texttt{background\_location\_} \texttt{service} & Native FFI bridge \\
\texttt{core/location\_engine/*} & FFI bindings (13 files) \\
\texttt{notification\_helper} & FCM platform code \\
\texttt{app\_theme.dart}, \texttt{app\_text\_theme.dart} & Style constants \\
\texttt{chart\_tokens.dart} & Style constants \\
\texttt{pune\_sample\_trail.dart} & Static fixture data \\
\texttt{app\_di.dart}, \texttt{app\_router.dart} & DI wiring + route table \\
\bottomrule
\end{tabular}
\end{table}

The exclusion removes 21 files and 4,154 lines. Raw coverage is 54.4\%
(16,895 of 31,039 instrumented lines); filtered coverage is 55.6\%
(16,595 of 29,824 lines). The gap between raw and filtered is small, but
the excluded volume is not: 13\% of instrumented source is structurally
untestable. We initially underestimated this. Our early exclusion list
covered only the thin wrappers (about 1,200 lines), and we treated 5\%
as a design-smell threshold. Adding the location engine's FFI bridge
layer pushed the exclusion past that threshold, which forced us to
re-examine whether the bridge code could be restructured. It could not:
the Dart bindings mirror the Rust API surface, and splitting them further
would add indirection without enabling testability.

The \texttt{Makefile} declares a \texttt{COVERAGE\_THRESHOLD} of 85\%.
This threshold applies to targeted module runs, not the full-codebase
metric. The full-codebase 55.6\% reflects modules not yet covered (screens,
some data services), not a policy that 55\% is acceptable.

The guiding principle: \emph{exclude honestly and report both numbers}.
Reporting only the filtered number hides the gap. Reporting only the raw
number penalises a codebase for having platform wrappers it cannot
eliminate. Both numbers together tell the real story.

% ----------------------------------------------------------------------------
\subsection{Mock-Boundary Architecture}
\label{sec:mocks}

Every platform dependency enters the testable codebase through one of two
gates:

\textbf{Gate 1: Constructor injection via Cubit dependencies.}
All business logic lives in Cubits (we enforce Cubit-only state management
via custom lint; see Section~\ref{sec:lints}). Each Cubit takes its
dependencies as constructor parameters. Tests instantiate the Cubit with
\texttt{mocktail} mocks:

\begin{lstlisting}[language=Java,caption={Typical Cubit test pattern}]
class MockAdminService extends Mock
    implements AdminService {}

blocTest<EmployeeShiftCubit, ShiftState>(
  'emits punched-in when shift starts',
  build: () => EmployeeShiftCubit(
    adminService: MockAdminService(),
    locationStorage: MockLocationStorage(),
  ),
  act: (c) => c.punchIn(officeId: 'office-1'),
  expect: () => [isA<ShiftActive>()],
);
\end{lstlisting}

This pattern scales. The test suite has 29 Cubit test files covering 28
of 50 production Cubits (56\% file-level coverage). The covered Cubits
are concentrated in the tracking and attendance domains, where bugs cost
the most; the 22 untested Cubits are admin, reporting, and announcement
screens where the business logic is thinner. The
\texttt{EmployeeShiftCubit} alone has 4 test files totalling 1,033 lines,
covering punch-in, punch-out, geofence checks, and session hydration.

\textbf{Gate 2: Platform interface stubs for Firebase.}
Firebase tests use \texttt{firebase\_core\_platform\_interface} (a dev
dependency) which provides \texttt{setupFirebaseCoreMocks()}. This
function registers a fake \texttt{FirebasePlatform} so that
\texttt{Firebase.initializeApp()} succeeds in the test process. On top of
this, individual Firebase services (\texttt{FirebaseAuth},
\texttt{FirebaseFirestore}) are mocked via \texttt{mocktail}.

One pattern deserves mention. The phone-auth wrapper
(\texttt{firebase\_phone\_auth.dart}, 233 lines) was initially excluded as
``untestable.'' We later wrote 390 lines of test code for it using a
callback-capture technique:

\begin{lstlisting}[language=Java,caption={Capturing Firebase verification callbacks}]
// Capture the callbacks Firebase calls async
when(() => mockAuth.verifyPhoneNumber(
  phoneNumber: any(named: 'phoneNumber'),
  verificationCompleted: captureAny(...),
  verificationFailed: captureAny(...),
  codeSent: captureAny(...),
  codeAutoRetrievalTimeout: captureAny(...),
)).thenAnswer((_) async {});

// Drive the captured callback synchronously
final codeSent = verify(...).captured.last;
codeSent('verification-id-123', null);
\end{lstlisting}

This file taught us that ``untestable'' often means ``expensive to test.''
The 390 lines of test scaffolding for 233 lines of production code is a
ratio of 1.67x. Whether that cost is worth it depends on the file's blast
radius. For phone authentication (the app's only login path), it was.

\textbf{GetIt reset pattern.}
A shared helper (\texttt{get\_it\_helper.dart}) calls
\texttt{GetIt.instance.reset()} and re-registers the service locator with
test mocks in each test group's \texttt{setUp}. This prevents
cross-contamination between tests that rely on different mock
configurations.

\textbf{SharedPreferences.}
Flutter's SDK provides \texttt{SharedPreferences.setMockInitialValues(\{\})}
for test use. We wrap this in a 7-line helper
(\texttt{app\_storage\_helper.dart}) that all storage tests import.

\textbf{Widget test theme scaffolding.}
The app uses a custom design system (\texttt{forui}) rather than raw
Material. Every widget test must pump its widget inside a theme provider,
localisations, and a \texttt{BlocProvider<ThemeCubit>}. We extracted this
into \texttt{wrapWithTheme()}, which has four variants: light, dark,
overlay-light, overlay-dark. Without this helper, each widget test file
would repeat 15--20 lines of ancestor setup. The helper is 82 lines and
is imported by every widget test in the suite. This is scaffolding that
the Flutter testing documentation does not warn you about: the moment you
leave Material defaults, every widget test needs a custom test harness.

\textbf{Driving futures inside \texttt{fakeAsync}.}
A 28-line helper (\texttt{fake\_async\_resolve.dart}) provides a generic
\texttt{resolve<T>()} function that drives a \texttt{Future<T>} to
completion inside \texttt{fakeAsync}, re-throwing any captured error. The
standard \texttt{fakeAsync} API does not provide this out of the box;
without it, Cubit methods that await a Future silently hang in the test
and produce a timeout with no useful stack trace.

% ----------------------------------------------------------------------------
\subsection{Custom Lint Rules as Architectural Enforcement}
\label{sec:lints}

The \texttt{zontra\_lints} package lives in \texttt{packages/zontra\_lints}
as a monorepo path dependency (published publicly as
\texttt{kx\_lints}\footnote{\url{https://pub.dev/packages/kx\_lints}}).
It uses the \texttt{custom\_lint\_builder}
framework (version 0.8.0) and the Dart \texttt{analyzer} API. It contains
17 rules at the time of writing. Table~\ref{tab:lints} lists the rules
relevant to testability.

\begin{table}[h]
\caption{Custom lint rules that enable testability}
\label{tab:lints}
\small
\begin{tabular}{@{}p{3.8cm}p{4.5cm}@{}}
\toprule
Rule & Effect on testability \\
\midrule
\texttt{cubit\_only} & Bans \texttt{extends Bloc<E,S>}; forces Cubit-only
  state management. Cubits are simpler to test: no event
  class boilerplate, direct method calls in \texttt{blocTest}. \\
\texttt{file\_length} & ERROR at 400 lines. Forces splitting large files,
  which produces smaller, more testable units. \\
\texttt{one\_widget\_per\_file} & One widget class per file (excludes
  \texttt{State<T>}). Keeps widget tests isolated. \\
\texttt{route\_constants\_only} & Navigation must use \texttt{AppRoutes.*}
  constants. Prevents string-literal routes that break silently. \\
\texttt{no\_raw\_alpha\_integer} & Born from PR \#421 (see
  Section~\ref{sec:regressions}). Bans integer alpha in
  \texttt{Color.withValues(alpha: N)} where N > 1.0. \\
\texttt{late\_controllers\_} \texttt{not\_final} & Prevents \texttt{late final} on
  animation/scroll controllers. Ensures \texttt{didUpdateWidget} can
  reassign. \\
\bottomrule
\end{tabular}
\end{table}

The remaining 11 rules enforce design-system consistency. Three examples:
\texttt{no\_material\_widgets} bans raw Material widgets
(\texttt{ElevatedButton}, \texttt{TextField}) in favour of project
wrappers that carry the design system's tokens.
\texttt{no\_theme\_of} bans \texttt{Theme.of(context)} and forces access
through typed token accessors, which eliminates a class of
\texttt{null} errors in widget tests where the theme context is
incomplete.
\texttt{widgets\_via\_barrel\_import} requires importing core widgets
through a barrel file, so renaming or moving a widget breaks one import
line instead of forty.

These rules do not directly test anything, but they shrink the surface
area of the UI layer. A widget test that pumps a project button does not
need to configure Material's full \texttt{ThemeData}; it configures the
project's thinner \texttt{FThemeData} instead.

Each rule is a single Dart class extending \texttt{DartLintRule} from
\texttt{custom\_lint\_builder}. The \texttt{no\_raw\_alpha\_integer} rule
(born from Case 1) is 47 lines: it visits every
\texttt{InstanceCreationExpression} and
\texttt{MethodInvocation} named \texttt{withValues}, checks whether the
\texttt{alpha} argument is a numeric literal greater than 1.0, and
reports a warning with a suggested fix. Writing a new rule takes two to
four hours, including tests for the rule itself.

Custom lint rules are cheaper than code review for architectural
constraints. A reviewer catches a violation once. A lint rule catches it
before the code reaches review, in every file, across every contributor,
permanently.
Sadowski et al.~\cite{sadowski2018lessons} reported similar findings on
the value of automated checks in Google's code review process.

% ============================================================================
% 5. REGRESSIONS CAUGHT
% ============================================================================

\section{Tests Catching Production Bugs}
\label{sec:regressions}

Six cases illustrate what automated testing bought us.

\textbf{Case 1: Alpha channel integer/float mixup (PR \#421).}
The \texttt{AppOtpField} widget passed 8-bit integers (160, 200, 120, 180)
to \texttt{Color.withValues(alpha:)}, which expects a 0.0--1.0 double. The
resulting colours were wrong on some devices, invisible on others. A widget
test was added after the fix. More importantly, the bug produced a new
custom lint rule (\texttt{no\_raw\_alpha\_integer}) that prevents the same
class of error codebase-wide. One bug became one test and one lint rule:
three layers of defence from a single incident.

\textbf{Case 2: Tracking data discrepancy (PR \#402).}
The employee view showed 22.9 km over 7 hours (423 waypoints). The admin
view showed 3.7 km in under a minute (11 of 62 waypoints). Five
compounding root causes: un-uploaded background waypoints, no flush on
reattach, wrong route reading from local SharedPreferences, ongoing
session zero-defaults, and UTC vs.\ IST time comparisons. Unit and widget
tests were added for the waypoint model, location storage, session trail,
and rider trail screen. The test count for the tracking module went from
12 to 47.

\textbf{Cases 3--4: Widget assertion fragility (commits \texttt{2641276},
\texttt{34394f8}).}
Two widget test failures taught us what to assert on. A greeting test
(``Good morning, \ldots'') passed in the morning and failed at night
because CI ran after hours; the fix was \texttt{textContaining}. A
separate refactor swapped \texttt{Text('No waypoints recorded today.')}
for a \texttt{NotFoundState} widget, and every test that pinned the old
copy broke; the fix was asserting on widget type instead of string
content. The rule we landed on: assert structure when it is stable, copy
when it is the contract.

\textbf{Case 5: Route-ordering security bypass (\texttt{app\_router\_test}).}
The attendance module defines two routes:
\texttt{/attendance/:userId/history} (static, admin-gated) and
\texttt{/attendance/:userId/:date} (parametric). Fluro matches
first-registered-wins. If history is registered after detail, both URLs
match the parametric route and the admin security gate in
\texttt{\_attendanceHistoryHandler} is silently bypassed. A route-ordering
test pins the registration sequence. This is the only case where a test
caught a security-relevant bug; the other five are correctness or
performance.

\textbf{Case 6: Test suite performance (commit \texttt{25865ef}, PR \#183).}
The \texttt{riders\_service\_test} took 15.3 seconds because the Cubit
under test used real timer delays. Conversion to \texttt{fakeAsync}
dropped it to under a second. \texttt{otp\_controller\_test} went from 6.4
seconds to near-instant. At the time of conversion (1,451 tests), the
full suite dropped from 34 to 22 seconds. The suite has since grown to
2,583 tests and takes longer, but the per-file speedups from
\texttt{fakeAsync} remain. This is not a production bug catch, but it
mattered for adoption: a slow test file is a test file developers skip.

% ============================================================================
% 6. COSTS AND NUMBERS
% ============================================================================

\section{Costs and Observations}
\label{sec:costs}

\begin{table}[h]
\caption{Test suite metrics}
\label{tab:metrics}
\small
\begin{tabular}{@{}lr@{}}
\toprule
Metric & Value \\
\midrule
Source files & 750 \\
Source lines (\texttt{lib/}) & 104,886 \\
Test files & 292 \\
Test invocations & 2,583 \\
Test lines & 53,354 \\
Test-to-source line ratio & 0.51 \\
Raw line coverage & 54.4\% (16,895 / 31,039) \\
Filtered line coverage & 55.6\% (16,595 / 29,824) \\
Excluded files & 21 (4,154 lines) \\
Coverage exclusion overhead & 13\% of instrumented source \\
Custom lint rules & 17 \\
\bottomrule
\end{tabular}
\end{table}

The test-to-source ratio of 0.51 means roughly one line of test for every
two lines of production code. This is lower than the 1:1 ratio sometimes
cited as a target, but higher than the median for mobile
apps~\cite{kochhar2015understanding}. The tracking and attendance Cubits
are closer to 1:1; the screen and widget layers pull the average down
because many screens lack widget tests.

Beyond test files themselves, the suite requires 18 helper files totalling
roughly 1,500 lines: shared mocks, fixture factories for domain models
(riders, tasks, teams, waypoints), a 136-line
\texttt{FirebaseRemoteConfigPlatform} fake that allows per-test feature
flag overrides, and a \texttt{flutter\_test\_config.dart} that installs a
tolerant golden comparator (0.5\% pixel diff threshold) and suppresses
two known Flutter 3.47 framework assertion regressions. This
infrastructure is invisible in test-to-source ratios but took weeks to
build and must be maintained when dependencies upgrade.

The framework regression suppression deserves a note. Flutter 3.47
introduced assertion failures in \texttt{MergeSemantics} (a forui
dependency) and \texttt{EditableText} deactivation that broke our suite
without any change on our side. The fix was a global test binding in
\texttt{flutter\_test\_config.dart} that catches and suppresses these
specific assertions. This is a maintenance cost that does not appear in
any metric: the suite would report hundreds of failures on a framework
upgrade, none of which reflect bugs in our code.

The test invocation count grew from roughly 1,450 to 2,583 over two
months. The growth came from two sources: golden tests for core design
system widgets (each widget generates multiple test invocations across
theme variants), and expanded coverage of the tracking module after the
data discrepancy bug (Case 2).

We have no generated code (\texttt{.g.dart}, \texttt{.freezed.dart}) in
the current codebase. This is unusual for a Flutter app of this size, and
it simplifies the exclusion policy. If we adopt code generation in the
future, the exclusion patterns already account for it. The project also
has six integration test files (1,355 lines) covering auth, offline sync,
push notifications, and shift lifecycle flows; these are outside the
scope of this paper but represent additional testing investment.

One cost we did not anticipate: test maintenance during UI refactors.
Widget tests that assert on specific text strings break on copy changes.
Widget tests that assert on widget types survive refactors but miss
regressions in the actual displayed content. We have not found a clean
answer. Our current rule: assert on type for layout structure, pin one
text string per screen as a smoke check.

% ============================================================================
% 7. LESSONS LEARNED
% ============================================================================

\section{Lessons Learned}
\label{sec:lessons}

\textbf{L1: ``Untestable'' usually means ``expensive to test.''}
We initially excluded the Firebase phone-auth wrapper as structurally
untestable. Then it caused two production bugs. We went back and wrote
390 lines of test scaffolding for 233 lines of production code, a 1.67x
ratio that would look wasteful on a spreadsheet.

It was not wasteful. That file is the only login path for every user of
the app. Cost-per-line is not the right metric; blast radius is.

\textbf{L2: One regression can fund an entire static-analysis rule.}
The alpha-channel bug (Case 1) took two hours to diagnose. The lint rule
took three hours to write. Every future developer on this codebase gets the
protection without knowing the bug existed. Custom lint rules are the
highest-leverage testing investment we made.

\textbf{L3: Coverage numbers need context or they mislead.}
55.6\% coverage sounds mediocre. But 13\% of our instrumented source is
platform code that cannot be unit-tested, and another ~20\% is screen
widgets we have not yet reached. The tracking and attendance Cubits,
where bugs actually hurt, are above 90\%. A single coverage number
without the exclusion
context is a vanity metric at best and a perverse incentive at worst.
Inozemtseva and Holmes~\cite{inozemtseva2014coverage} showed that coverage
alone is not strongly correlated with fault detection; our experience
confirms this.

\textbf{L4: Lint rules beat code review for architectural invariants.}
A reviewer might miss a \texttt{Bloc<E,S>} in a PR with 400 changed lines.
The \texttt{cubit\_only} lint rule will not. We enforce Cubit-only state
management, single-widget-per-file, 400-line file caps, barrel imports,
and route constants through lint rules rather than review checklists. The
rules fire in the IDE before the developer opens a PR.

\textbf{L5: Fast test files matter more than fast suites.}
Converting \texttt{riders\_service\_test} from real timers to
\texttt{fakeAsync} dropped it from 15.3 seconds to under one. That
single change made the difference between running the file after every
edit and running it once before pushing. The full suite has since grown
past the point where running it after saving is practical (2,583 tests),
but per-file speed still determines whether a developer tests at all
during a focused editing session.

% ============================================================================
% 8. THREATS TO VALIDITY
% ============================================================================

\section{Threats to Validity}
\label{sec:threats}

\textbf{External validity.} This is a single-project experience report.
Zontra's dependency profile (Firebase, FFI, SharedPreferences) is common
but not universal; apps using other backends (Supabase, AWS Amplify,
custom REST) will face different testing boundaries. The exclusion
framework transfers; the specific exclusion patterns may not.

\textbf{Internal validity.} We attribute regression catches to the test
suite, but we cannot prove that manual testing would not have caught
the same bugs. Case 2 (the tracking discrepancy) was first reported by a
user, not by a test. The tests were written \emph{after} the fix to
prevent recurrence, not as a proactive catch.

\textbf{Construct validity.} We report line coverage as our primary
metric. Line coverage is an imperfect proxy for test
effectiveness~\cite{inozemtseva2014coverage,gopinath2014code}. We do not
claim that 55.6\% coverage means the app is 55.6\% ``tested''; we use it
as a directional measure and a ratchet to prevent regression of tested
modules.

\textbf{Researcher bias.} The first author is also the primary developer
of the codebase and the test suite. This dual role means the testing
decisions and their evaluation are not independent. We mitigate this by
reporting concrete metrics (line counts, coverage percentages, commit
hashes) that a reviewer can verify against the repository, but we cannot
rule out selection bias in which regression cases we chose to highlight.

\textbf{Team size.} The development team is two engineers. A larger team
might face coordination costs (conflicting mocks, flaky shared fixtures)
that we do not encounter. The custom lint rules may become more valuable
at larger team sizes, but we cannot verify this.

% ============================================================================
% 9. RELATED WORK
% ============================================================================

\section{Related Work}
\label{sec:related}

\textbf{Mobile testing.}
Kochhar et al.~\cite{kochhar2015understanding} surveyed 83 Android
developers and found that 14\% wrote no tests at all;
Linares-V\'{a}squez et al.~\cite{linares2017developers} found a median
test ratio of 0.04 across 1,000 Android apps. Our ratio of 0.51 is an
order of magnitude higher, which we attribute partly to the Cubit
architecture making business logic injection-friendly by default. On the
challenges side, Muccini et al.~\cite{muccini2012software} identified
platform fragmentation and Coppola et
al.~\cite{coppola2019mobile} studied GUI test fragility; our widget test
breakages (Cases 3--4) match their findings. Zein et
al.~\cite{zein2016systematic} noted that most research targets native
Android/iOS; cross-platform testing, especially for Flutter, remains
under-studied. We found no peer-reviewed work on unit-testing Flutter
apps with Firebase or FFI dependencies. The closest practitioner
literature is Flutter's own testing documentation~\cite{flutter_testing_docs},
which covers the framework's API but not the scaffolding required when
production code depends on platform services that the test harness
cannot initialise.

\textbf{Cross-platform development.}
Bi{\o}rn-Hansen et al.~\cite{biorn2020flutter} surveyed cross-platform
approaches and noted that testing support varies significantly across
frameworks. Flutter's testing layer is more mature than most
alternatives, but the gap between documentation and production practice
that we describe has not been previously reported.

\textbf{Code review and static analysis.}
Sadowski et al.~\cite{sadowski2018lessons} reported on code review at
Google, finding that automated checks reduced review burden and caught
issues that human reviewers missed. Our custom lint rules serve a
similar function at a much smaller scale.
Beller et al.~\cite{beller2019developer} studied developer testing
behaviour in the IDE and found that developers who receive faster
feedback write more tests, consistent with our observation (L5) about
suite speed.

\textbf{Coverage effectiveness.}
Inozemtseva and Holmes~\cite{inozemtseva2014coverage} found that
coverage is not strongly correlated with test suite effectiveness after
controlling for suite size. Gopinath et al.~\cite{gopinath2014code}
qualified this: coverage is useful as a relative
measure within a project but misleading as an absolute quality target.
Our dual-number reporting (raw and filtered) is an attempt to make
coverage useful without making it a target that distorts behaviour.

% ============================================================================
% 10. CONCLUSION
% ============================================================================

\section{Conclusion}
\label{sec:conclusion}

The exclusion policy, mock boundaries, and lint rules took us from zero
tests to 2,583 in a codebase where 13\% of the instrumented source is
structurally untestable. But 55.6\% filtered coverage is a waypoint, not
a destination, and the hardest open problem is one we have not solved:
widget tests that are specific enough to catch regressions but stable
enough to survive UI refactors without constant maintenance. We do not
know whether our current heuristic (assert type for structure, pin one
text string per screen) will hold as the app grows past 150,000 lines.

What we are more confident about: custom lint rules returned more per
hour than anything else in this effort. A rule costs hours to write and
catches every future instance of a bug class, across every file, without
a reviewer needing to remember it. If a team can afford only one thing
from this report, it should be that.

% ============================================================================
% REFERENCES
% ============================================================================

\bibliographystyle{IEEEtran}
\bibliography{references}

\end{document}
